% !TeX root = ../main.tex
\documentclass[../main.tex]{subfiles}

\begin{document}
\chapter*{Introducción}
\addcontentsline{toc}{chapter}{Introducción}

El presente apunte se organiza por \textbf{bloques temáticos} y no por una secuencia estricta de clases. El objetivo de esta decisión editorial es mantener una exposición más cercana a la estructura conceptual del curso y, al mismo tiempo, conservar una base suficientemente flexible como para incorporar desarrollos, ejemplos y comentarios sin desordenar el conjunto.

En particular, la redacción de esta versión busca mantenerse lo más próxima posible al estilo de las fuentes de referencia del curso, especialmente a los archivos \texttt{electrodinamica.tex} y \texttt{cap-Electrodinamica.tex}. Por ello, cuando se añade algún comentario interpretativo o una aclaración física, se procura hacerlo sin alterar las ecuaciones ni el hilo lógico de las demostraciones.

\begin{cajaidea}
\textbf{Criterio editorial seguido en este apunte:}
\begin{enumerate}[label=\arabic*.]
    \item organizar el contenido según la \textbf{dependencia conceptual} entre los temas;
    \item conservar, dentro de cada bloque, el orden lógico en que las ideas aparecen en el curso;
    \item distinguir con claridad entre \textbf{desarrollo principal}, \textbf{observaciones} y \textbf{notas de cuidado};
    \item evitar reformulaciones que modifiquen el contenido matemático de las demostraciones originales.
\end{enumerate}
\end{cajaidea}

En esta primera parte, el material queda concentrado principalmente en la estructura local de la electrodinámica clásica, la obtención de las ecuaciones de onda electromagnéticas en el vacío, el formalismo de potenciales y la construcción causal de los potenciales retardados.

\newpage

\end{document}
